\section*{Bab \ref{ch:g9:ions} --- Ion dan Larutan Ion}

\begin{solution}{exo:g9:ions:1}
\omterm{def:g7:atoms-and-molecules:atom}{Atom} atau sekelompok \omterm{def:g7:atoms-and-molecules:atom}{atom} yang telah melepaskan atau menerima \omterm{def:g9:inside-the-atom:nucleus}{elektron}
dan bermuatan. \omterm{def:g9:ions:ion}{Kation} bermuatan positif (\omterm{def:g9:inside-the-atom:nucleus}{elektron} dilepaskan), \omterm{def:g9:ions:ion}{anion}
negatif (\omterm{def:g9:inside-the-atom:nucleus}{elektron} diterima).
\end{solution}

\begin{solution}{exo:g9:ions:2}
\ce{Mg^{2+}}: 12 \omterm{def:g9:inside-the-atom:proton}{proton}, 10 \omterm{def:g9:inside-the-atom:nucleus}{elektron}.
\end{solution}

\begin{solution}{exo:g9:ions:3}
\ce{KCl}, \ce{MgCl2}, \ce{CaO}.
\end{solution}

\begin{solution}{exo:g9:ions:4}
Air garam mengandung \omterm{def:g9:ions:ion}{ion}, \ce{Na+(aq)} dan \ce{Cl-(aq)}, yang bergerak
dan membawa muatan. \omterm{def:g7:atoms-and-molecules:molecule}{Molekul} gula tidak bermuatan.
\end{solution}

\begin{solution}{exo:g9:ions:5}
\ce{CuSO4}, \ce{AlCl3}, \ce{Na2CO3}.
\end{solution}

\begin{solution}{exo:g9:ions:6}
Besi(III), \ce{Fe^{3+}}: \ce{Fe^{3+}(aq) + 3OH-(aq) -> Fe(OH)3(s)}.
\end{solution}

\begin{solution}{exo:g9:ions:7}
Sulfat: kelebihan 2 \omterm{def:g9:inside-the-atom:nucleus}{elektron}. Nitrat: 1.
\end{solution}

\begin{solution}{exo:g9:ions:8}
Putih dengan perak nitrat: \ce{Cl-}. Biru dengan natrium hidroksida:
\ce{Cu^{2+}}. Tembaga(II) klorida, \ce{CuCl2}.
\end{solution}

\begin{solution}{exo:g9:ions:9}
Empat: satu di setiap sisi (kiri, kanan, atas, bawah).
\end{solution}

\begin{solution}{exo:g9:ions:10}
\ce{Fe^{2+}(aq) + 2OH-(aq) -> Fe(OH)2(s)}.
\end{solution}

\begin{solution}{exo:g9:ions:11}
Seng klorida: \omterm{def:g9:ions:ion}{ion} natrium tidak memberikan \omterm{def:g9:ions:precipitate}{endapan} dengan natrium
hidroksida, \omterm{def:g9:ions:ion}{ion} seng memberikan \omterm{def:g9:ions:precipitate}{endapan} putih:
\ce{Zn^{2+}(aq) + 2OH-(aq) -> Zn(OH)2(s)}.
\end{solution}

\begin{solution}{exo:g9:ions:12}
\ce{CaCl2}: 2000 \omterm{def:g9:ions:ion}{ion} klorida. Muatan: $1000 \times (+2) + 2000 \times
(-1) = 0$.
\end{solution}

\begin{solution}{pb:g9:ions:1}
\textbf{1.} Besi(II), \ce{Fe^{2+}}: \ce{Fe(OH)2}.

\textbf{2.} Besi(III), \ce{Fe^{3+}}: \ce{Fe(OH)3}.

\textbf{3.} \ce{Fe^{2+}} memiliki satu lebih banyak: $26 - 2 = 24$
\omterm{def:g9:inside-the-atom:nucleus}{elektron}, dibandingkan $26 - 3 = 23$ untuk \ce{Fe^{3+}}.

\textbf{4.} \ce{Fe^{3+}(aq) + 3OH-(aq) -> Fe(OH)3(s)}.

\textbf{5.} Air itu keluar mengandung \omterm{def:g9:ions:ion}{ion} besi(II); ketika bersentuhan
dengan \omterm{def:g4:air-a-mixture-of-gases:air}{udara}, \omterm{def:g9:ions:ion}{ion-ion} itu menjadi \omterm{def:g9:ions:ion}{ion} besi(III), yang membentuk padatan
jingga.

\textbf{6.} \ce{Fe(OH)3}: $(+3) + 3 \times (-1) = 0$.

\textbf{7.} Tidak ketika baru keluar: \omterm{def:g9:ions:ion}{ion-ion} besi terlarut dan lolos
melewati saringan. Begitu padatan jingga terbentuk, ya: padatan itu
\omterm{def:g9:ions:precipitate}{endapan}, yang tertahan oleh saringan.

\textbf{8.} Bukan: air itu \omterm{def:g6:pure-substances-and-mixtures:mixture}{campuran} (air dan \omterm{def:g9:ions:ion}{ion-ion} terlarut). Ketika
baru diambil dan masih jernih, air itu homogen.

\textbf{9.} $56 \times 1.67 \times 10^{-27} = \qty{9.35e-26}{kg}$.

\textbf{10.} $\qty{0.30}{mg} = \qty{0.30e-6}{kg} = \qty{3.0e-7}{kg}$.

\textbf{11.} Sebuah \omterm{def:g9:ions:ion}{ion} berbeda dari \omterm{def:g7:atoms-and-molecules:atom}{atomnya} hanya dua atau tiga
\omterm{def:g9:inside-the-atom:nucleus}{elektron}, yang massanya dapat diabaikan dibandingkan intinya.

\textbf{12.} $3.0 \times 10^{-7} / 9.35 \times 10^{-26} \approx 3.2
\times 10^{18}$ \omterm{def:g9:ions:ion}{ion} besi per liter.
\end{solution}
